\section{When linear beats gradient boosting --- the TSO-feature flip}

The most quotable finding in the repo. Same features, same walk-forward,
same two years (\texttt{reports/backtests/2026-07-16\_2yr\_summary.csv}
and the 12-month campaign):

\begin{center}
\begin{tabular}{lcc}
setup & ridge & LightGBM \\
\hline
without the TSO forecast & 4.05\% & 3.16\% \quad (GBM wins by 0.9pp) \\
with the TSO forecast    & \textbf{2.08\%} & 2.12\% \quad (ridge edges it) \\
\end{tabular}
\end{center}

Gradient boosting is not overrated --- adding one feature CHANGED THE
PROBLEM, and the new problem favors a different model class.

\subsection{Three mechanisms}

\begin{enumerate}
  \item \textbf{The task became residual correction.} With PSE's
        forecast as input, the model no longer forecasts load; it
        corrects the small error of an expert who already modeled the
        nonlinearities (temperature response, holidays, ramps). What
        remains is nearly linear. Trees' edge IS nonlinearity; there is
        none left to harvest. Sensitivity analysis agrees: the TSO
        feature carries 96\% of total skill; all else fights over
        $\sim$0.1pp.
  \item \textbf{Variance rules near the noise floor.} TSO alone: 2.23\%.
        Irreducible noise is not far below. In that regime the winner is
        whoever adds least estimation variance. Ridge: $\sim$25 shrunk
        coefficients from 8{,}760 rows. GBM: hundreds of leaf values,
        each from a data subset. More variance, no bias advantage left
        to buy with it.
  \item \textbf{Feature engineering pre-linearized the problem.}
        hour sin/cos makes daily seasonality linear-fittable;
        heating/cooling degree-days linearize the V-shaped temperature
        response. Good features + linear model beats raw features +
        flexible model --- the same reason LEAR (linear) embarrasses deep
        nets across the price-forecasting literature.
\end{enumerate}

\subsection{The control experiment we already ran}

Price is the counter-case: spiky, regime-switching, genuinely nonlinear
target with no expert forecast to lean on --- and there LightGBM DOES
beat the linear model (MAE 17.8 vs 18.5). The rule that survives both
products:

\begin{quote}
Trees win where nonlinear structure survives in the target. Linear wins
where an expert forecast or your own feature engineering already
removed it.
\end{quote}

Interview line: ``Gradient boosting beat ridge by 0.9pp until both got
the TSO's forecast --- then ridge won, because the leftover problem is
linear residual correction and the lowest-variance model takes it.
Knowing when the fancy model stops paying is the skill.''
